\section{The slope is the homology of the smoothing closure}
\label{slope-section}

The finite-threshold recurrence identifies its slope as the homology dimension
of a repeated differential. We now identify that number topologically. This
also proves that the repeated part cannot disappear. All vector spaces in this
section are over $\F$, and the reduction point is fixed on a closing arc
outside the run boxes. Let
\[
 V=\bigoplus_q E^q,\qquad F=D_E+L,\qquad
 L=X_p+X_q,\qquad
 \eta=\dim V-2\operatorname{rank}F,
\]
where $p$ and $q$ lie on the cap and cup of the distinguished $E$ smoothing.
Here $D_E$ is the context differential. The notation $X_p$ denotes the
usual dot action, namely multiplication by $X$ on the resolution circle
containing $p$. Let $J_E$ be the actual link diagram obtained by replacing
the entire distinguished run by that smoothing. The identification below
concerns total homology dimension; it does not identify the original
homological grading on the summands of $V$ with a grading on $H(V,F)$.

\subsection{The two dots lie on one actual component}

\begin{lemma}\label{slope-cap-cup}
The cap and cup of the distinguished smoothing belong to the same component
of $J_E$, for every choice of the other signed braid runs.
\end{lemma}

\begin{proof}
Remove the interior of the smoothing box and retain the four cut ends. The
outside diagram consists of closed components and two open paths. The braid
orientation implies that each open path joins a top cut end to a bottom
cut end: two ends of the same kind cannot be the ends of a coherently
oriented outside path. Consequently the outside matching of these four ends is
one of the two top-to-bottom matchings. Joining the two top ends by the cap
and the two bottom ends by the cup produces a single closed component in
either case. This component contains both smoothing arcs. The argument
concerns the actual strands, with crossings retained. In a complete
resolution the two arcs can lie on different circles, so it would be
incorrect to set $L=0$ state by state.
\end{proof}

\subsection{The characteristic-two weight slide}

Jaeger's weighted-dot construction proves that, in characteristic two,
relocating markings while preserving their total weight on every actual
component preserves the twisted complex~\cite[Theorem~3.1]{jaeger}. We give
the local algebra needed here, including its compatibility with reduction.
It supplies an explicit change of basis, a stronger statement than merely
saying that two dot actions induce the same homology map.

\begin{lemma}\label{slope-slide}
Let $(K,d)$ be the ordinary crossing-cube complex of a link, with its
gradings forgotten. Suppose $p,p'$ are successive points of one actual
strand, separated by one crossing. If $W$ is any sum of fixed-arc dot
actions and $W'=W+X_p+X_{p'}$, then $(K,d+W)$ and $(K,d+W')$ are isomorphic
as vector spaces equipped with square-zero endomorphisms. The isomorphism
preserves the reduced subcomplex at a fixed marked point.
\end{lemma}

\begin{proof}
Use the Frobenius algebra $A=\F[X]/(X^2)$, with multiplication $m$ and
comultiplication
\[
 \Delta(1)=1\otimes X+X\otimes1,\qquad
 \Delta(X)=X\otimes X.
\]
Direct evaluation on the basis gives
\begin{align}
 m\Delta&=0,&
 \Delta m&=X\otimes1+1\otimes X,\label{slope-frobenius}\\
 m(X\otimes1)&=Xm=m(1\otimes X),&
 (X\otimes1)\Delta&=\Delta X=(1\otimes X)\Delta.
 \label{slope-naturality}
\end{align}
The notation on the right of the first line means the corresponding
endomorphism of $A\otimes A$.
Dot operators commute and square to zero; they also commute with the cube
differential by \eqref{slope-naturality}. Thus $(d+W)^2=0$ for every $W$
in the statement.

Let $H$ be the reverse saddle at the selected crossing: it changes its cube
coordinate from $1$ to $0$, uses the saddle opposite to the ordinary
differential, and is zero in the other direction. Thus $H^2=0$.
Saddles at other cube coordinates commute with $H$; their two occurrences
in $dH+Hd$ cancel in characteristic two. At the selected crossing,
\eqref{slope-frobenius} gives
\[
 dH+Hd=X_p+X_{p'}=:T.
\]
Moreover $TH=0$. If the reverse saddle merges, its two output dot actions
agree and sum to zero. If it splits, the assertion follows from
$(X\otimes1+1\otimes X)\Delta=0$. Equation~\eqref{slope-naturality}
shows that each fixed-arc dot commutes with $H$, including the dots
appearing in $W$. Therefore
\begin{align*}
 (d+W')(1+H)+(1+H)(d+W)
 &=dH+Hd+(W'+W)+W'H+HW\\
 &=T+T+WH+HW+TH=0.
\end{align*}
Since $(1+H)^2=1$, this proves the claimed isomorphism.

If $z$ is the reduction point, then $H$ commutes with $X_z$ by the same
local identities. Hence $1+H$ preserves $\operatorname{im}X_z$, the
subcomplex in which the marked resolution circle is labeled $X$. Its
inverse does so as well. This also covers the case where $z$ belongs to a
different component from the moving dots.
\end{proof}

Repeatedly applying the lemma moves a dot along any path in its actual
component. Two unit weights meeting at one point cancel because $1+1=0$.
By Lemma~\ref{slope-cap-cup}, the reduced cube of $J_E$ with differential
$d+X_p+X_q$ is consequently isomorphic to its ordinary reduced cube.
This argument uses a sequence of local invertible maps. A general
null-homotopy $T=dH+Hd$ would not, by itself, justify replacing $d+T$ by
$d$; the identities $H^2=0$ and $TH=0$ are what make the displayed
conjugation work.

\subsection{Compatibility with the compressed context}

The preceding argument applies initially to the crossing cube. Passing to
the run complex requires compatibility of the compression with the two
exterior dots. This follows from the local cobordism construction, but it
is useful to record exactly which homotopy identities are needed.

\begin{lemma}\label{slope-compression}
Suppose $(K,d)$ and $(V,D_E)$ are related by the local Bar--Natan
reductions in the context run boxes. Let $L_K$ and $L_V$ be the sums of the
two dot actions on the distinguished smoothing arcs. Then these reductions
induce a homotopy equivalence between $(K,d+L_K)$ and $(V,D_E+L_V)$ after
forgetting the gradings. They also preserve the fixed reduction.
\end{lemma}

\begin{proof}
Bar--Natan's tangle formalism and its delooping and cancellation operations
take place in the category of local dotted cobordisms~\cite{bn05,bn07}.
An exterior marked arc is retained throughout each such operation. Sliding
a dot along its unchanged boundary cylinder shows that every local
structural map, and every local contracting homotopy, commutes with that
dot. This remains true for linear combinations and composites, and for
the maps obtained by delooping a circle inside a context box. The two
distinguished dots and the reduction point are outside those boxes.

More explicitly, choose the resulting maps $f:K\to V$, $g:V\to K$ and
homotopies $h$ on $K$, $k$ on $V$. In addition to their ordinary chain-map
identities they satisfy
\[
 fL_K=L_Vf,\quad gL_V=L_Kg,\quad
 hL_K=L_Kh,\quad kL_V=L_Vk.
\]
Thus $f$ and $g$ are chain maps for the deformed differentials. Their
homotopy identities remain valid, because
\[
 gf-1=dh+hd=(d+L_K)h+h(d+L_K),
\]
and similarly
\[
 fg-1=D_Ek+kD_E=(D_E+L_V)k+k(D_E+L_V).
\]
The additional terms cancel in characteristic two. All maps commute with
the reduction dot, so the same identities restrict to the reduced
subcomplexes. Only finite complexes are involved; no completion or
infinite-twist limit enters this argument.
\end{proof}

\begin{theorem}[Topological slope]\label{slope-theorem}
For either sign of the distinguished braid run,
\begin{equation}\label{slope-identity}
 \boxed{\displaystyle
 \eta=\dim_{\F}\rKh(J_E;\F)>0.}
\end{equation}
The reduction point on the right is the original fixed point, and total
dimension is taken over all homological and quantum degrees.
\end{theorem}

\begin{proof}
Lemma~\ref{slope-compression} identifies $H(V,F)$ with the homology of the
reduced crossing cube deformed by the two dots. The weight-slide argument
identifies that cube with the ordinary reduced cube of $J_E$. This proves
the equality and shows that the slope is independent of the sign chosen
for the distinguished run.

For strict positivity, the graded Euler characteristic of ordinary reduced
integral Khovanov homology is the reduced Jones polynomial, up to the
standard normalization~\cite{bn05}. For a nonempty link with $\ell$
components, its value at $1$ has absolute value $2^{\ell-1}$ and is
nonzero. Thus the integral reduced homology has a nonzero free summand.
The universal coefficient theorem then gives nonzero reduced homology
over $\F$. The diagram $J_E$ is nonempty, so this applies even when it is
a link with several components.
\end{proof}

\begin{remark}\label{slope-grading}
The dot deformation combines maps of different original bidegrees.
Jaeger's formulation retains a collapsed $\delta$ grading, and forgetting
that grading is sufficient for \eqref{slope-identity}. The formula does
not recover the individual homological or quantum degrees of $\rKh(J_E)$
inside the repeated part. The homological-degree profile for the original
twisted closure follows instead from the separate identification of its
successive complete $E$ groups and both incident boundary maps. These are
distinct assertions and use different grading information.
\end{remark}

\subsection{Computational consequences and output size}

Write the context interval as $[a,c]$, let $w=c-a$, and put $m_0=w+2$.
For $m\ge m_0$, set $\delta=m-m_0$. The two signed recurrence formulas
have exactly $\delta$ inserted degrees. Combining them with
Theorem~\ref{slope-theorem} gives
\begin{equation}\label{slope-rank-growth}
 R(m)=R(m_0)+\delta\,\dim_{\F}\rKh(J_E;\F),
\end{equation}
where $R$ denotes total reduced rank for either fixed choice of sign.
Every inserted degree is nonzero. In particular, this is an exact linear
growth statement with positive integer slope, not merely an upper bound
on the possible growth.

The identity provides an independent interpretation and a possible
cross-check of the slope. The algorithm need not compute $\rKh(J_E)$
separately: one reference computation at $m_0$ already contains a central
boundary equal to $F$, so that $\eta=\dim V-2\operatorname{rank}F$.
The two reference shoulder Betti numbers need not equal $\eta$; the
boundary-rank extraction is essential. The original closure's component
count remains necessary before using the Khovanov unknot-detection
criterion~\cite{km}, because shortening an exponent can change the
permutation parity and the number of components.

There is also an unconditional output-size obstruction. An explicit list
of nonzero degree--Betti-number pairs contains at least $\delta$ entries,
and therefore requires $\Omega(\delta)$ output operations. For fixed
context and growing binary-encoded $m$, this is exponential in the bit
length of $m$. A compact profile avoids that cost by retaining the two
fixed shoulders, their degree shifts, and one constant-coefficient
interval. Its dependence on the distinguished exponent is only through
binary integers of $O(\log m)$ bits. The same is true of the extra
arithmetic needed to return \eqref{slope-rank-growth}. This observation
does not bound the cost of computing the context data as the context
grows; that cost can still be exponential. It establishes the appropriate
output representation and the precise sense in which a single large
exponent can be handled without expansion.

\subsection{Position relative to existing tangle methods}

The local twist normal form, dot calculus, and cancellation machinery are
established mathematics~\cite{bn05,bn07}. Thompson gives compact zig-zag
representatives for rational-tangle Khovanov complexes~\cite{thompson}.
Jaeger's weight-moving theorem supplies the central algebra behind the
slope identification~\cite{jaeger}. Kotelskiy, Watson, and Zibrowius give
immersed-curve descriptions of four-ended tangle invariants, gluing
theorems recovering Khovanov homology, and an interpretation of the
two-strand infinite-twist construction~\cite{kwz}. These precedents make
clear why a general priority claim for twist stability would be
inappropriate.

The contribution developed here is the explicit finite threshold, the
extraction from one reference computation, the compact exact output, and
their implementation and verification relative to the pinned ProveIt
baseline. The immersed-curve viewpoint is a plausible route for studying
several variable twist regions, but obtaining useful complexity bounds
would require control of the context representation, local systems, and
gluing calculations. The one-run theorem alone supplies no such control
and does not justify a general quasi-polynomial unknot-recognition bound.
